%% kapa: exam=1 round=21 year=2010 period=2 form=A subject=accounting subject_ko=회계학 seq=13 number=13
%% source: Q-Net 감정평가사 기출문제 (cst00309) artlSeq=5024181 file=extracted/1008/5024181/2010년 제21회 감정평가사 1차시험 기출문제/2교시.hwp
%% conversion: kordoc 4.24.51 -> tools/gen1.py (auto); review=proofread-gpt-sol; flags=none
\begin{question}{13}
확정기여제도와 확정급여제도에 관한 설명으로 옳지 않은 것은?
\begin{choices}
  \item 확정기여제도에서는 기업이 별개의 실체에 고정 기여금을 납부한다.
  \item 확정급여제도에서 기업의 의무는 약정한 퇴직급여를 종업원에게 지급하는 것이다.
  \item 확정기여제도보다 확정급여제도의 회계처리가 비교적 복잡하다.
  \item 확정급여제도에서는 채무와 비용의 측정에 보험수리적 가정이 요구된다.
  \item 확정기여제도에서는 기업이 적립금의 투자위험을 부담한다.
\end{choices}
\end{question}
\answer{5}
